% ============================================================================
% OPERATOR UI ARCHITECTURE
% ============================================================================

\section{Operator UI}
\label{section:ui-architecture}

\subsection{Executive summary}

The STAR operator UI is a browser-based cockpit written in TypeScript on
top of React 19 and bundled with Vite. The runtime is a single-page
application: the operator opens the Pi 5 in a browser, the page
fetches its bundle from the static-asset server, and the client opens a
single WebSocket to the Go gateway (\S\ref{section:repository-layout}).
All telemetry and command traffic flows through that one WebSocket as
binary-encoded \texttt{STAREnvelope} messages -- the same Protocol
Buffer wire contract as the firmware (\S\ref{section:nanopb-protocol})
but multiplexed at the UI level into a single envelope per direction.
A Zustand store fans incoming envelopes out to per-panel selectors so a
new lidar scan does not re-render the battery panel. Layout is a tiling
window manager built on \texttt{react-grid-layout} with persisted
per-view overrides. Three panels (\texttt{Nav2GoalPanel},
\texttt{PidTuningPanel}, \texttt{TransportDiagPanel}) are explicitly
placeholder UIs: they render a yellow "Not implemented yet" banner and
disable their action buttons, because the gateway does not yet expose
the corresponding RPC over the UI WebSocket. They exist so the layout
slots are reserved; they do not fake state.

\subsection{Layout overview}

The UI is a flat tree under \texttt{star-ui/src/}:

\begin{center}
\begin{tabular}{|p{3.6cm}|p{11.2cm}|}
\hline
\textbf{Path}                    & \textbf{Role}                                                                                                                              \\
\hline
\texttt{main.tsx}                & React entry point. Mounts \texttt{App} inside a custom \texttt{ErrorBoundary} that prints uncaught exceptions to a monospace fallback page. \\
\texttt{App.tsx}                 & Top-level component. Resolves the active view, builds the panel registry, runs the responsive grid, owns the maximize-overlay UX.            \\
\texttt{theme.ts}                & A single \texttt{COLORS} record of semantic accent / status colors plus theme-aware CSS variable references.                                  \\
\texttt{components/}             & 17 panels + the status bar + the joystick \texttt{ControllerView}. Each panel is self-contained and reads from one or two store selectors.   \\
\texttt{hooks/}                  & \texttt{useGamepad} (Gamepad API polling) and \texttt{useSTARConnection} (WebSocket lifecycle + reconnect + REST fallback for E-Stop).        \\
\texttt{store/}                  & Two Zustand stores: \texttt{dashboardStore} (telemetry, alerts, e-stop) and \texttt{useWindowStore} (active view, per-view layout overrides). \\
\texttt{services/GatewayService.ts} & Packet ring buffer + envelope-preview formatter for the packet analyzer.                                                                  \\
\texttt{proto/star/v1/}          & Generated TypeScript bindings (\texttt{protobuf-ts}) for every \texttt{star.v1} message used on the wire.                                    \\
\texttt{test/setup.ts}, \texttt{hooks/useGamepad.test.ts} & Vitest setup + the only unit test currently shipped (gamepad polling, deadzone, axes mapping).                          \\
\hline
\end{tabular}
\end{center}

The data flow inside the browser:

\begin{center}
\begin{verbatim}
   Gateway (Go on Pi5)
     |
     |  binary STAREnvelope (Protocol Buffers)
     v
   WebSocket (ws://<pi-host>:8080/ws)
     |
     v
   useSTARConnection                  navigator.getGamepads()
     |  STAREnvelope.fromBinary           |
     |                                    v
     |                                useGamepad
     |                                    |  { linearVel, angularVel }
     v                                    v
   dashboardStore                     TeleopPanel @ 50 Hz
     | telemetry / motors /                |  ControllerState envelope
     | battery / odometry /                v
     | lidar / system /                useSTARConnection.sendControllerState
     | alerts / eStop                      |  STAREnvelope.toBinary
     v                                     v
   useShallow / .subscribe()            WebSocket (TX path)
     |
     v
   Panels (MotorPanel, BatteryPanel, OdometryPanel, ...)
     |
     |  StatusBar -> sendEStop / sendEStopRelease
     |   (REST /api/estop fallback if WS is down)
     |
   useWindowStore (zustand + persist via localStorage)
     |
     v
   ResponsiveGridLayout (react-grid-layout)
\end{verbatim}
\end{center}

\subsection{Panels}

Every entry under \texttt{star-ui/src/components/} is one row below.
The \emph{Status} column distinguishes three categories. \emph{Live}
means the panel renders real gateway data and has no stub banner.
\emph{Read-only} means the panel only displays data and never writes
back. \emph{Placeholder (banner)} means the panel renders a yellow
"Not implemented yet" banner and disables its action buttons because
the corresponding gateway RPC is not yet wired to the UI WebSocket.

\begin{center}
\begin{tabular}{|p{2.6cm}|p{4.7cm}|p{2.6cm}|p{4.0cm}|}
\hline
\textbf{Panel}            & \textbf{Purpose}                                                            & \textbf{Status}             & \textbf{Data source}                                                       \\
\hline
\texttt{TeleopPanel}      & Joystick / gamepad teleop. 50 Hz \texttt{ControllerState} TX.               & Live                        & \texttt{useGamepad} + \texttt{sendControllerState}                          \\
\texttt{ControllerView}   & Visual joystick + connection pill rendered inside \texttt{TeleopPanel}.    & Live (read-only)            & Props from \texttt{TeleopPanel}                                            \\
\texttt{MotorPanel}       & Per-wheel velocity bars (FL / FR / BL / BR).                                & Live (read-only)            & \texttt{dashboardStore.motors}                                              \\
\texttt{BatteryPanel}     & Pack voltage, SOC, temperature, segmented battery icon.                     & Live (read-only)            & \texttt{dashboardStore.battery}                                             \\
\texttt{OdometryPanel}    & 2D odometry trail on a canvas plus x / y / heading readouts.                & Live (read-only)            & \texttt{dashboardStore.odometry}                                            \\
\texttt{LidarPanel}       & Polar LiDAR scatter on a canvas; \texttt{store.subscribe} (no re-render).   & Live (read-only)            & \texttt{dashboardStore.lidarScan}                                            \\
\texttt{ImuPanel}         & Roll / pitch / yaw attitude indicator + accel + gyro readouts.              & Live (read-only)            & \texttt{dashboardStore.telemetry.imu}                                       \\
\texttt{GpsPanel}         & Lat / lon / alt / accuracy + fix-quality badge.                             & Live (read-only)            & \texttt{dashboardStore.telemetry.gps}                                       \\
\texttt{SystemHealthPanel}& Uptime, free heap, CPU, temperature, firmware version, WiFi RSSI.           & Live (read-only)            & \texttt{dashboardStore.systemStatus} + \texttt{telemetry}                   \\
\texttt{AlertsPanel}      & Severity-colored alert cards with relative age.                             & Live (read-only)            & \texttt{dashboardStore.alerts}                                              \\
\texttt{DiagLogPanel}     & Monospace scrolling log of every alert (newest first).                      & Live (read-only)            & \texttt{dashboardStore.alerts}                                              \\
\texttt{TimeSeriesPanel}  & 1 Hz sparklines: FL / FR motor velocity, CPU \%, temperature.               & Live (read-only)            & \texttt{dashboardStore} (polled at 1 Hz via getState)                       \\
\texttt{PacketAnalyzer}   & Last 30 envelopes (TX + RX) with seq / type / size / decoded preview.       & Live (read-only)            & \texttt{services/GatewayService.getRecentPackets()} ring buffer (200 ms poll) \\
\texttt{StatusBar}        & Connection pill + STALE / robot mode / E-Stop / RESUME / Layouts / theme.   & Live                        & \texttt{dashboardStore} + \texttt{sendEStop} + \texttt{sendEStopRelease}     \\
\texttt{CameraFeed}       & Camera placeholder (icon + reticle); no WebRTC / MJPEG client yet.          & Placeholder (no banner)     & None -- static SVG only                                                     \\
\texttt{FirmwarePanel}    & OTA UI shell with simulated progress bar (\texttt{simulateUpload}).         & Placeholder (no banner)     & Local state only -- no gateway RPC wired                                    \\
\texttt{Nav2GoalPanel}    & Inputs for X / Y / theta navigation goals; Send button disabled.            & Placeholder (banner)        & None -- gateway has no Nav2 goal RPC                                        \\
\texttt{PidTuningPanel}   & Kp / Ki / Kd inputs per motor (FL / FR / BL / BR); Apply button disabled.   & Placeholder (banner)        & None -- no UI-side channel to \texttt{SetMotorPidConfig}                    \\
\texttt{TransportDiagPanel}& Static SPI / USB CDC health rows; no live data.                            & Placeholder (banner)        & None -- gateway does not stream \texttt{TransportDiagnostics} on \texttt{/ws} \\
\hline
\end{tabular}
\end{center}

\noindent\textbf{Note on placeholders:} \texttt{CameraFeed} and
\texttt{FirmwarePanel} are placeholder UIs but do not render the
"banner" pattern -- the camera shows a "No Camera Stream" empty state,
and the firmware panel exposes interactive buttons that drive a local
simulated progress timer. The three banner-bearing panels
(\texttt{Nav2GoalPanel}, \texttt{PidTuningPanel},
\texttt{TransportDiagPanel}) are the ones honest about being stubs.

\subsection{State management}

Two Zustand stores own all UI state. Both are kept deliberately small
and per-slice; panels select only the field they need via
\texttt{useShallow} or a stable getter so a new envelope re-renders only
the relevant subtree.

\subsubsection{\texttt{store/dashboardStore.ts}}

The telemetry feed. It owns connection state, the last seq number, the
"is the data stale" flag, and one slice per major payload kind. The
critical method is \texttt{updateFromEnvelope(env)}: each \texttt{case}
in the oneof switch performs exactly one \texttt{set()} call so React
re-renders only the slice that changed. A sequence-gap detector raises
a \texttt{SEQ\_GAP} \texttt{WARN} alert when \texttt{seq > lastSeq + 1}.

\begin{center}
\begin{tabular}{|p{3.6cm}|p{4.0cm}|p{6.7cm}|}
\hline
\textbf{Field}              & \textbf{Type}                            & \textbf{Source}                                                          \\
\hline
\texttt{connectionState}    & \texttt{ConnectionState} (4-state)       & \texttt{useSTARConnection} on open / close.                              \\
\texttt{lastSeq}            & \texttt{number}                          & \texttt{env.seq} for gateway-originated envelopes.                       \\
\texttt{seqGapDetected}     & \texttt{boolean}                         & Set when a gap is observed (for the Status bar).                          \\
\texttt{dataIsStale}        & \texttt{boolean}                         & Set on \texttt{ws.onclose}; cleared on first \texttt{telemetry} envelope. \\
\texttt{telemetry}          & \texttt{TelemetryData|null}              & oneof case \texttt{telemetry}.                                            \\
\texttt{motors}             & \texttt{MotorStatus[]|null} (length 4)   & oneof case \texttt{motors}.                                               \\
\texttt{battery}            & \texttt{BatteryState|null}               & oneof case \texttt{battery}.                                              \\
\texttt{odometry}           & \texttt{OdometryData|null}               & oneof case \texttt{odometry}.                                             \\
\texttt{lidarScan}          & \texttt{LidarScan|null}                  & oneof case \texttt{lidar}.                                                \\
\texttt{systemStatus}       & \texttt{SystemStatus|null}               & oneof case \texttt{system}.                                               \\
\texttt{alerts}             & \texttt{Alert[]} (capped at 200)         & oneof case \texttt{alert} + locally-injected \texttt{SEQ\_GAP} entries.   \\
\texttt{eStopActive}        & \texttt{boolean}                         & \texttt{triggerEStop} / \texttt{releaseEStop} from the StatusBar.         \\
\hline
\end{tabular}
\end{center}

\subsubsection{\texttt{store/useWindowStore.ts}}

The layout / view manager. Six named views (\texttt{OVERVIEW},
\texttt{TELEOP}, \texttt{NAVIGATION}, \texttt{DEBUG}, \texttt{CONFIG},
\texttt{FULL}); each declares its panel set and a default
\texttt{react-grid-layout} for the \texttt{lg} breakpoint. The
\texttt{FULL} view is the only scrollable one. The store is wrapped in
\texttt{persist(...)} so per-view drag / resize edits survive a reload
under the localStorage key \texttt{robot-dashboard-v6}. Reset clears the
overrides for the active view, not all views.

\begin{center}
\begin{tabular}{|p{3.4cm}|p{4.3cm}|p{6.6cm}|}
\hline
\textbf{Field}              & \textbf{Type}                                  & \textbf{Purpose}                                                  \\
\hline
\texttt{activeView}         & \texttt{ViewName}                              & One of \texttt{OVERVIEW} / \texttt{TELEOP} / \texttt{NAVIGATION} / \texttt{DEBUG} / \texttt{CONFIG} / \texttt{FULL}. \\
\texttt{viewLayouts}        & \texttt{Partial<Record<ViewName,Layouts>>}     & Per-view user overrides; absent keys fall back to the view's default. \\
\texttt{maximizedPanel}     & \texttt{string|null}                           & The key of the currently maximized panel (Esc closes).            \\
\texttt{setView}            & action                                         & Switch view; clears \texttt{maximizedPanel}.                      \\
\texttt{updateLayouts}      & action                                         & Persist a drag / resize for the current view.                     \\
\texttt{toggleMaximize}     & action                                         & Open / close the maximize overlay.                                \\
\texttt{resetLayout}        & action                                         & Drop overrides for the active view only.                          \\
\hline
\end{tabular}
\end{center}

\subsection{Custom hooks}

\begin{center}
\begin{tabular}{|p{3.4cm}|p{11.4cm}|}
\hline
\textbf{Hook}                  & \textbf{Behaviour}                                                                                                                        \\
\hline
\texttt{useGamepad}            & Polls \texttt{navigator.getGamepads()} on every animation frame. Picks the first non-null gamepad, applies a 0.1 deadzone to both axes, returns \texttt{\{ connected, linearVel, angularVel \}} where \texttt{linearVel = -axes[1]} (forward stick = +1) and \texttt{angularVel = axes[0]}. Disconnect events fall back to the polling tick rather than firing immediately. \\
\texttt{useSTARConnection}     & Owns the lifetime of a single \texttt{WebSocket}. Sets \texttt{binaryType = 'arraybuffer'} \emph{before} any messages arrive (the browser default \texttt{blob} would crash \texttt{Uint8Array} construction). On every message it calls \texttt{STAREnvelope.fromBinary}, then \texttt{updateFromEnvelope} on the store, then \texttt{recordPacket} into the analyzer ring. On close it marks data stale and reconnects with exponential backoff (\texttt{500\,ms} base, \texttt{30\,s} cap, \texttt{$\pm$30\,\%} jitter). Returns three TX functions: \texttt{sendControllerState(state)}, \texttt{sendEStop(reason)}, \texttt{sendEStopRelease()}. The two e-stop senders fall back to a REST \texttt{POST /api/estop?action=...\&reason=...} if the WebSocket is not open -- this is the safety-critical guarantee that operator e-stop never silently fails. \\
\hline
\end{tabular}
\end{center}

The repository ships exactly these two custom hooks (plus the
\texttt{useGamepad.test.ts} Vitest suite). All other state lives in
\texttt{useDashboardStore} or \texttt{useWindowStore}.

\subsection{Service layer}

\texttt{services/GatewayService.ts} is intentionally narrow. It does
\emph{not} own the WebSocket, the encoder, or the decoder. It owns one
data structure: a \texttt{RING\_SIZE = 100} entry ring buffer of
\texttt{PacketRecord}s, written by \texttt{recordPacket(env, dir,
sizeBytes, errorType?)} and read by \texttt{getRecentPackets()}. The
ring is module-scoped (not in Zustand) so that high-rate envelope
arrivals never trigger a store re-render -- the
\texttt{PacketAnalyzer} polls the buffer at 200\,ms intervals instead.
Each record is \texttt{\{ seq, type, direction, sizeBytes, tsMs,
preview \}}, and \texttt{formatPreview} pretty-prints the oneof
variant for the \texttt{DECODED} column (e.g. for a \texttt{lidar}
payload it renders \texttt{"<n> pts"}; for \texttt{controller} it
renders \texttt{"lin=...\,ang=...\,[TX]"}). The file also exports
\texttt{MV\_PER\_V} as the canonical \texttt{milliVolt -> Volt}
divisor used by \texttt{BatteryPanel} so battery math is consistent
across the UI.

The actual gRPC-over-WebSocket plumbing -- \texttt{STAREnvelope.toBinary}
on TX and \texttt{STAREnvelope.fromBinary} on RX -- lives in
\texttt{useSTARConnection}, not here. The \texttt{*.client.ts} files
under \texttt{src/proto/star/v1/} (\texttt{configuration.client.ts},
\texttt{firmware\_update.client.ts}, \texttt{motor\_control.client.ts},
\texttt{telemetry.client.ts}) are protobuf-ts gRPC clients that the UI
\emph{generates} but does not yet instantiate. The intent is that
configuration / OTA / PID-set RPCs will be served via gRPC-Web through
a separate transport once those panels are de-stubbed; today the only
live transport is the multiplexed \texttt{STAREnvelope} on
\texttt{/ws}.

\subsection{WebSocket protocol}

The UI side speaks one wire format: \texttt{star.v1.STAREnvelope},
encoded as Protocol Buffers binary (no length prefix -- WebSocket
already frames). The envelope carries metadata
(\texttt{seq}, \texttt{ts\_ms}) plus a \texttt{oneof payload}.

\begin{center}
\begin{tabular}{|p{2.6cm}|p{3.4cm}|p{8.6cm}|}
\hline
\textbf{Direction} & \textbf{Sender}    & \textbf{Payload variants used}                                                                                                  \\
\hline
RX (gateway -> UI) & Gateway            & \texttt{telemetry}, \texttt{motors}, \texttt{battery}, \texttt{odometry}, \texttt{lidar}, \texttt{system}, \texttt{alert}.       \\
TX (UI -> gateway) & UI                 & \texttt{controller} (50 Hz from \texttt{TeleopPanel}), \texttt{estop} (operator E-Stop / Resume from \texttt{StatusBar}).        \\
\hline
\end{tabular}
\end{center}

\noindent\textbf{Latency expectations:} the TX cadence is 50\,Hz
(20\,ms period, \texttt{sendIntervalMs = 1000 / sendRateHz} in
\texttt{TeleopPanel.tsx}). The RX cadence is whatever the gateway
chooses to push; per-payload rates are governed gateway-side. The UI
makes no rate assumption on RX -- the \texttt{seq} field plus the
\texttt{seqGapDetected} flag are the only freshness signals it
exposes. On \texttt{ws.onclose} the store is marked stale immediately;
panels render their last-known value but the \texttt{StatusBar} shows
the \texttt{STALE} pill so the operator does not trust it.

\noindent\textbf{Out-of-band fallback:} the only TX path with a
non-WebSocket fallback is operator E-Stop, which posts to
\texttt{/api/estop} via \texttt{fetch()} when the WebSocket is closed.
This is the dual-channel safety-critical path required by
\S\ref{section:safety-architecture}.

\subsection{Build and run}

\begin{center}
\begin{tabular}{|p{3.4cm}|p{11.4cm}|}
\hline
\textbf{Command}                & \textbf{What it does}                                                                                                                       \\
\hline
\texttt{npm install}            & Install dependencies (React 19, Zustand 5, react-grid-layout, protobuf-ts, uPlot, Vite 7, Vitest 4).                                          \\
\texttt{npm run dev}            & Vite dev server on \texttt{0.0.0.0:5173} (\texttt{strictPort: true}). HMR enabled.                                                            \\
\texttt{npm run build}          & \texttt{tsc -b \&\& vite build}. Strict TypeScript build first, then Vite production bundle into \texttt{dist/}.                              \\
\texttt{npm run preview}        & Serve the built \texttt{dist/} bundle (production smoke test).                                                                                \\
\texttt{npm run lint}           & ESLint over the whole project (\texttt{eslint .}).                                                                                            \\
\texttt{npm run test}           & \texttt{vitest run} -- single-shot Vitest run with the \texttt{jsdom} environment.                                                            \\
\hline
\end{tabular}
\end{center}

\noindent\textbf{Environment variables:} the only env var the runtime
reads is \texttt{VITE\_WS\_PORT} (parsed in \texttt{App.tsx} line 31).
If absent, the UI defaults to port 8080. Protocol (\texttt{ws://} vs
\texttt{wss://}) is auto-selected from \texttt{window.location.protocol}.
Hostname always tracks \texttt{window.location.hostname}, so the same
bundle works whether the operator opens the Pi directly
(\texttt{http://star.local}) or through Caddy
(\texttt{https://star.local}). There is no \texttt{VITE\_STAR\_WS\_URL}
override -- the URL is constructed from the three components above.

\subsection{Theming and layout}

\subsubsection{Theme}

\texttt{src/theme.ts} exports a \texttt{COLORS} record. Static accent
colors (\texttt{accent}, \texttt{warning}, \texttt{success},
\texttt{danger}, \texttt{estopActive}) are literal hex values and stay
the same in both themes -- they are SF / iOS system colors. Theme-aware
colors (\texttt{textPrimary}, \texttt{panelBg}, \texttt{border}, etc.)
are not literals at all, they are the strings
\texttt{var(--color-text)}, \texttt{var(--panel-bg)}, etc., so they
resolve at runtime through CSS custom properties defined in
\texttt{index.css}. The \texttt{StatusBar} writes
\texttt{document.documentElement.dataset.theme = 'dark'|'light'} and
persists the choice to \texttt{localStorage['star-theme']}.

\subsubsection{Layout engine}

\texttt{App.tsx} renders one \texttt{<ResponsiveGridLayout>} from
\texttt{react-grid-layout}. Five breakpoints
(\texttt{lg}/\texttt{md}/\texttt{sm}/\texttt{xs}/\texttt{xxs} with
column counts \texttt{24}/\texttt{18}/\texttt{12}/\texttt{8}/\texttt{6})
let panels reflow on tablets and phones. Row height for non-scrollable
views is recomputed by a \texttt{ResizeObserver} so 18 rows always fill
the viewport minus the 80\,px top bar and the 12\,px gaps. Drag is
restricted to \texttt{.panel-header}; resize handles are exposed on all
eight edges. A maximize button on each panel opens a fullscreen
overlay (Esc closes). All panel chrome (the glass background, the
hover, the stagger-fade-in animation) lives in \texttt{App.css} /
\texttt{index.css}, not in TypeScript.

\subsection{Testing}

The Vitest setup is intentionally light. \texttt{src/test/setup.ts}
imports \texttt{@testing-library/jest-dom} for matchers; the
environment is \texttt{jsdom} (declared in \texttt{vitest.config.ts}).
The repository ships exactly one test suite,
\texttt{src/hooks/useGamepad.test.ts}, covering the three gamepad
behaviours (initial disconnected state, axes-to-velocity mapping,
0.1 deadzone). \texttt{npm run test} runs the suite once and exits.
There are no Playwright / Cypress browser tests today; UI-side
end-to-end coverage is delegated to the gateway-level e2e harness in
\texttt{star-gateway/test/e2e/}.

\subsection{Where to look in code}

\begin{center}
\begin{tabular}{|p{6.0cm}|p{8.7cm}|}
\hline
\textbf{If you want to...}                                  & \textbf{Open}                                                                                                       \\
\hline
Add a new panel                                              & \texttt{star-ui/src/components/<Name>Panel.tsx}, then register the key in \texttt{App.tsx} \texttt{usePanelRegistry} and add a layout entry to one or more views in \texttt{store/useWindowStore.ts}. \\
Add a new RX payload kind                                    & Extend the \texttt{oneof} in \texttt{star-proto/proto/star/v1/ui.proto}, regen, then add a \texttt{case} branch in \texttt{store/dashboardStore.ts} \texttt{updateFromEnvelope}. \\
Change WebSocket URL / port                                  & \texttt{App.tsx} (the \texttt{WS\_PORT} / \texttt{WS\_PROTOCOL} / \texttt{WS\_URL} block) and / or set \texttt{VITE\_WS\_PORT}. \\
Tune the gamepad deadzone or rate                            & \texttt{src/hooks/useGamepad.ts} (deadzone constant) and \texttt{src/components/TeleopPanel.tsx} (\texttt{sendRateHz}). \\
De-stub a placeholder panel                                  & \texttt{Nav2GoalPanel.tsx} / \texttt{PidTuningPanel.tsx} / \texttt{TransportDiagPanel.tsx} -- remove the banner and wire the disabled button to the gateway RPC. \\
Add a new view (panel preset)                                & \texttt{store/useWindowStore.ts} -- add a \texttt{ViewName} entry to \texttt{views} with its panel list and \texttt{lg} layout. The \texttt{StatusBar} dropdown picks them up automatically. \\
Reset persisted layouts                                      & In the browser devtools, clear \texttt{localStorage['robot-dashboard-v6']}; or click "Reset to Default" in the Layouts dropdown for the active view only. \\
Inspect TX / RX traffic                                      & Open the \texttt{DEBUG} view -- the \texttt{PacketAnalyzer} panel shows the last 30 envelopes with type, size, age and a decoded preview. \\
Add a unit test                                              & \texttt{src/hooks/<thing>.test.ts}; \texttt{npm run test}. The \texttt{useGamepad} suite is the template. \\
\hline
\end{tabular}
\end{center}
